(1) 一次透射振幅乘积是 \(t_1t_2\)，每次额外往返贡献 \(r_1r_2e^{i\phi}\)。取镜面振幅约定使 \(|r_j|=\sqrt{R_j}\)、\(|t_j|^2=1-R_j\)，求和后 \(E_t/E_i=t_1t_2e^{i\phi/2}/(1-r_1r_2e^{i\phi})\)。因此 \(\boxed{\mathcal T(\phi)=\frac{(1-R_1)(1-R_2)}{1+R_1R_2-2\sqrt{R_1R_2}\cos\phi}}\)。

(2) 在 \(\phi=2j\pi\) 的峰值 \(\mathcal T_{\max}=(1-R_1)(1-R_2)/(1-\sqrt{R_1R_2})^2\)。因为 \((1-\sqrt{R_1R_2})^2-(1-R_1)(1-R_2)=(\sqrt{R_1}-\sqrt{R_2})^2\ge0\)，故峰值不超过1，恰当且仅当 \(R_1=R_2\) 时完全透射。

(3) 自由光谱范围为 \(\Delta\nu_{\rm FSR}=c/(2nL)\)。记 \(\rho=\sqrt{R_1R_2}\)，峰附近分母 \(\simeq(1-\rho)^2+\rho(\phi-2j\pi)^2\)，所以往返相位的FWHM约为 \(2(1-\rho)/\sqrt\rho\)。精细度 \(\boxed{\mathcal F\simeq\pi\sqrt\rho/(1-\rho)}\)，频率线宽 \(\Delta\nu=\Delta\nu_{\rm FSR}/\mathcal F\)。

(4) 实际吸收或散射在几何级数中给附加往返振幅因子，因而降低峰透射并增宽线宽。共振条件 \(2nL\nu/c=j\)，若 \(n\) 近似固定，\(\delta\nu/\nu=-\delta L/L\)；若折射率也漂移，则 \(\delta\nu/\nu=-\delta(nL)/(nL)\)。
